% GENERATED FILE. Edit canonical lessons, then run npm run generate:books.
% canonical-source-hash: fnv1a64:f9453df7e24daaea
% canonical-lessons: JA-C85-utsu, JA-C85-suu, JA-C85-maku, JA-C85-taku, JA-C85-haku

\chapter{To hit, To breathe in, To sow, To cook, To put on}
\label{ch:ja-to-hit-to-breathe-in-to-sow-to-cook-to-put-on}
\hlchaptermodality{85}

\begin{chapteropening}
\textbf{By the end of this chapter:} \emph{I can say to hit, to breathe in, to sow, to cook and to put on.}

Complete the last lesson of chapter 85: i can say to hit, to breathe in, to sow, to cook and to put on.
\end{chapteropening}

\section[\texorpdfstring{\ja{うつ} (utsu)}{うつ (utsu)}]{\ja{うつ} (utsu) — to hit}

\label{lesson:JA-C85-utsu}



\begin{quote}
Before the new one: say the Japanese for this sort of, then the Japanese for that sort of.
\end{quote}

\subsection*{You'll want to know: \ja{うつ}}
\textbf{\ja{うつ}} — \emph{utsu} — \textquotedblleft{}to hit\textquotedblright{}.

\begin{grammarlens}[title={What you've built}]
One more word for this chapter.
\end{grammarlens}

\subsection*{Your turn}
\begin{itemize}
\raggedright
  \item \emph{Say it:} \emph{utsu}
  \item \emph{Say it:} \emph{utsu}, once more
  \item \emph{Say it:} say \emph{utsu}
  \item \emph{Recall:} say the Japanese for this sort of, then the Japanese for that sort of, then say \emph{utsu} again
\end{itemize}

\begin{tcolorbox}[breakable,colback=teal!4,colframe=teal!35!black,title={Before you move on}]
What does \ja{うつ} mean? (\textquotedblleft{}To hit\textquotedblright{}.) Say it once more.
\end{tcolorbox}
\section[\texorpdfstring{\ja{すう} (suu)}{すう (suu)}]{\ja{すう} (suu) — to breathe in}

\label{lesson:JA-C85-suu}



\begin{quote}
Before the new one: say the Japanese for that sort of, then the Japanese for to hit.
\end{quote}

\subsection*{You'll want to know: \ja{すう}}
\textbf{\ja{すう}} — \emph{suu} — \textquotedblleft{}to breathe in\textquotedblright{}.

\begin{grammarlens}[title={What you've built}]
One more word for this chapter.
\end{grammarlens}

\subsection*{Your turn}
\begin{itemize}
\raggedright
  \item \emph{Say it:} \emph{suu}
  \item \emph{Say it:} \emph{suu}, once more
  \item \emph{Say it:} say \emph{suu}
  \item \emph{Recall:} say the Japanese for that sort of, then the Japanese for to hit, then say \emph{suu} again
\end{itemize}

\begin{tcolorbox}[breakable,colback=teal!4,colframe=teal!35!black,title={Before you move on}]
What does \ja{すう} mean? (\textquotedblleft{}To breathe in\textquotedblright{}.) Say it once more.
\end{tcolorbox}
\section[\texorpdfstring{\ja{まく} (maku)}{まく (maku)}]{\ja{まく} (maku) — to sow; to wind}

\label{lesson:JA-C85-maku}



\begin{quote}
Before the new one: say the Japanese for to hit, then the Japanese for to breathe in.
\end{quote}

\subsection*{You'll want to know: \ja{まく}}
\textbf{\ja{まく}} — \emph{maku} — \textquotedblleft{}to sow; to wind\textquotedblright{}.

\begin{grammarlens}[title={What you've built}]
One more word for this chapter.
\end{grammarlens}

\subsection*{Your turn}
\begin{itemize}
\raggedright
  \item \emph{Say it:} \emph{maku}
  \item \emph{Say it:} \emph{maku}, once more
  \item \emph{Say it:} say \emph{maku}
  \item \emph{Recall:} say the Japanese for to hit, then the Japanese for to breathe in, then say \emph{maku} again
\end{itemize}

\begin{tcolorbox}[breakable,colback=teal!4,colframe=teal!35!black,title={Before you move on}]
What does \ja{まく} mean? (\textquotedblleft{}To sow; to wind\textquotedblright{}.) Say it once more.
\end{tcolorbox}
\section[\texorpdfstring{\ja{たく} (taku)}{たく (taku)}]{\ja{たく} (taku) — to cook (rice)}

\label{lesson:JA-C85-taku}



\begin{quote}
Before the new one: say the Japanese for to breathe in, then the Japanese for to sow; to wind.
\end{quote}

\subsection*{You'll want to know: \ja{たく}}
\textbf{\ja{たく}} — \emph{taku} — \textquotedblleft{}to cook (rice)\textquotedblright{}.

\begin{grammarlens}[title={What you've built}]
One more word for this chapter.
\end{grammarlens}

\subsection*{Your turn}
\begin{itemize}
\raggedright
  \item \emph{Say it:} \emph{taku}
  \item \emph{Say it:} \emph{taku}, once more
  \item \emph{Say it:} say \emph{taku}
  \item \emph{Recall:} say the Japanese for to breathe in, then the Japanese for to sow; to wind, then say \emph{taku} again
\end{itemize}

\begin{tcolorbox}[breakable,colback=teal!4,colframe=teal!35!black,title={Before you move on}]
What does \ja{たく} mean? (\textquotedblleft{}To cook (rice)\textquotedblright{}.) Say it once more.
\end{tcolorbox}
\section[\texorpdfstring{\ja{はく} (haku)}{はく (haku)}]{\ja{はく} (haku) — to put on (shoes, trousers)}

\label{lesson:JA-C85-haku}



\begin{quote}
Before the new one: say the Japanese for to sow; to wind, then the Japanese for to cook.
\end{quote}

\subsection*{You'll want to know: \ja{はく}}
\textbf{\ja{はく}} — \emph{haku} — \textquotedblleft{}to put on (shoes, trousers)\textquotedblright{}.

\begin{grammarlens}[title={What you've built}]
One more word for this chapter.
\end{grammarlens}

\subsection*{Your turn}
\begin{itemize}
\raggedright
  \item \emph{Say it:} \emph{haku}
  \item \emph{Say it:} \emph{haku}, once more
  \item \emph{Say it:} say \emph{haku}
  \item \emph{Recall:} say the Japanese for to sow; to wind, then the Japanese for to cook, then say \emph{haku} again
\end{itemize}

\begin{tcolorbox}[breakable,colback=teal!4,colframe=teal!35!black,title={Before you move on}]
What does \ja{はく} mean? (\textquotedblleft{}To put on (shoes, trousers)\textquotedblright{}.) Say it once more.
\end{tcolorbox}
